% ==============================================================================
% ssec:cs_ksvs_zemedelstvi --- Případová studie: KSVS zemědělství a výživa 2026
% Zdroje: KSVS_OSPZV_ASO_2026, sdeleni_47_2026, sdeleni_17_2026,
%         CMKOS_ZpravaKV2025, CMKOS_PrinosKV2025, zakon_kv_1991, zakon_zp_2006,
%         eu_directive_2022_2041
% ==============================================================================

na~rok 2026, uzavřená na rok 2026 mezi \acs{ZS ČR}, \acs{ČMSZP} a~odborovým svazem \acs{OSPZV}--\acs{ASO}, byla podepsána 19.~prosince 2025 v~Praze. Smlouva byla uložena na~\acloc{MPSV} dle \S~9 \acgen{ZKV} a~uložení bylo oznámeno Sdělením č.~17/2026~Sb. Její závaznost byla následně pro rok 2026 rozšířena Sdělením č.~47/2026~Sb.~na~všechny zaměstnavatele v~odvětvích \acgen{NACE}~01.1 až~01.6. Jde o~příklad použití mechanismu extenze podle \S~7 \acgen{ZKV}.\par

Extenze závaznosti chrání signatářské podniky sdružené v~zaměstnavatelských svazech před vnitroodvětvovým externím mzdovým dumpingem ze strany neorganizovaných subjektů. V~kontextu kap.~\ref{sec:teorie} ilustruje tato dohoda funkční odvětvovou strukturu mzdového a~pracovního standardu. Přesah této \acgen{KSVS} se ovšem přes zapojení stejného \acs{OS} nevztahuje na~maloobchod, Penny Market spadající pod \acs{NACE}~47.11.\par

Sektorová mzdová úprava definuje dvanáct stupňů minimálních mzdových tarifů pro stanovenou týdenní pracovní dobu \SI{40}{\hour\per\week}, s~rozpětím od~\SI{22400}{\czk\per\month} v~prvním stupni (kopírujícím zákonné minimum pro~rok 2026) do~\SI{45150}{\czk\per\month} ve~dvanáctém stupni. Tyto tarify jsou nepodkročitelné a~ostatní zaměstnavatelé v~rámci rozšířené platnosti je musí respektovat bez ohledu na~konkrétní interní uspořádání mzdové stupnice, což sjednocuje minima napříč odvětvím.\par

Sjednané mzdové příplatky překračují zákonné standardy \acgen{ZP}. Příplatek za~přesčasovou práci v~noční době nebo ve~dnech nepřetržitého odpočinku v~týdnu činí \SI{40}{\percent} \acgen{PV} (standardně \SI{25}{\percent}), příplatek za~běžnou práci v~noci \SI{20}{\percent} \acgen{PV} (\SI{10}{\percent}) a~za~práci v~sobotu a~v~neděli rovněž \SI{20}{\percent} \acgen{PV} (\SI{10}{\percent}).\par

Organizace pracovní doby zohledňuje charakter práce v~zemědělství. Vyrovnávací období pro nerovnoměrné rozvržení činí \SI{52}{\week}, přičemž sezónní období nižší potřeby práce je vymezeno od~1.~prosince do~konce února. Smlouva povoluje flexibilitu v~dodatečném dorovnání zkráceného nepřetržitého odpočinku mezi směnami v~rámci následujících tří týdnů a~týdenního odpočinku v~rámci tří týdnů v~úhrnu \SI{105}{\hour}, příp.~sezónně v~rámci šesti týdnů v~celkovém úhrnu \SI{210}{\hour}, což umožňuje vykrýt sklizňové špičky.\par

Základní dovolená činí \SI{4}{\week}, avšak je doplněna bonusem k~náhradě mzdy ve~výši nejméně \SI{2}{\percent} průměrného výdělku za~období od ledna do~listopadu. Tento finanční nárok zaniká, pokud podniková \ac{KS} nebo vnitřní předpis sjednává přímo pátý týden dovolené -- smluvní mechanismus tak stimuluje zavedení pátého týdne volna, resp.~kompenzuje jeho absenci. Sociální podpora zahrnuje měsíční příspěvek na~penzijní připojištění a~doplňkové spoření odstupňovaný podle délky \acgen{PP} u~zaměstnavatele: \SI{1000}{\czk\per\month} po odpracování jednoho roku, \SI{1100}{\czk\per\month} po dvou letech a~\SI{1350}{\czk\per\month} po třech a~více letech. Smlouva obsahuje též inflační doložku zavazující k~novému vyjednávání o~tarifech při růstu cenového indexu podle \acgen{ČSÚ} o~více než \SI{3}{\percent}.\par

Tato \ac{KSVS} tak poskytuje funkční kompromis: sezónní flexibilitu pracovní doby a~eliminaci dumpingu pro zaměstnavatele výměnou za transparentní strukturu odměňování a~zvýšený příplatkový standard pro zaměstnance. Nicméně plný dopad extenze v~odvětví zůstává podle analýz \acgen{ČMKOS} limitován úzkým věcným rozsahem rozšiřujícího sdělení \acloc{MPSV}.~\cite{KSVS_OSPZV_ASO_2026}~\cite{CMKOS_ZpravaKV2025}\par
